\chapter{Descripción general}
\label{cap:descripcion}

\section{El sistema}
\label{sec:elsistema}

\textbf{BeeTracer} es una plataforma de trazabilidad logística desarrollada por \textbf{Bilix
Ingeniería} (Valparaíso, 2015) que digitaliza el ciclo completo de una operación de recepción y
despacho de carga en terminales extraportuarios y bodegas. Su propuesta de valor central es la
\textbf{evidencia fotográfica obligatoria}: el sistema no permite avanzar de una etapa a la
siguiente sin haber capturado el registro visual correspondiente, de modo que al cerrar la faena se
dispone automáticamente de un expediente completo de lo ocurrido ---\textbf{quién, cuándo y
cómo}---.

El sistema se comercializa bajo un modelo de \textbf{licencia por suscripción} (mensual,
trimestral o anual) y se encuentra operativo en Chile, con prospectos en Perú, Argentina y Brasil.
Dado que el formato XML del manifiesto es un estándar mundial, la internacionalización requiere
únicamente adaptar la dirección del \este{web service} de la aduana local.

\section{Arquitectura general}
\label{sec:arquitectura}

BeeTracer se estructura en \textbf{dos frentes que se sincronizan en tiempo real}, descritos en la
Tabla~\ref{tab:arquitectura}.

{\estilotabla
\begin{longtable}{L{0.17\textwidth} L{0.20\textwidth} L{0.21\textwidth} L{0.33\textwidth}}

\caption{Componentes de la arquitectura general de BeeTracer}
\label{tab:arquitectura}
\\

\toprule
\textbf{Componente} & \textbf{Tecnología} & \textbf{Usuarios} & \textbf{Función} \\
\midrule
\endfirsthead

\toprule
\textbf{Componente} & \textbf{Tecnología} & \textbf{Usuarios} & \textbf{Función} \\
\midrule
\endhead

\textbf{Plataforma Web} &
Node.js (versión original en Java) &
Administrador, Jefe de faena, Supervisor &
Centro de control, configuración, planificación, revisión y análisis. Operada desde oficina. \\

\midrule

\textbf{Aplicación Móvil} &
Ionic (multiplataforma) &
Jefe tarjador / operario de terreno &
Captura de datos y evidencia fotográfica en patio. Diseñada para uso rudo sobre tablet o
celular. \\

\midrule

\textbf{Servidores} &
Servicio externo contratado &
--- &
Alojamiento de aplicación, base de datos y repositorio de imágenes. SLA de resolución estimado en
30 minutos ante caídas. \\

\midrule

\textbf{Seguridad} &
Certificados SSL, encriptación &
--- &
Protección de las transacciones y del acceso. \\

\bottomrule

\end{longtable}
}

\begin{nota}
\textbf{Nota histórica.} La primera versión del sistema se construyó en \textbf{Java}, con
\textbf{JasperReports} para la generación de los PDF y \textbf{SQL Server} como motor de base de
datos (exigencia del primer gran cliente). El núcleo tardó entre 3 y 4 meses en desarrollarse. La
versión moderna migró a \textbf{Node.js} en la web y a \textbf{Ionic} en el móvil.
\end{nota}

\section{Módulos funcionales}
\label{sec:modulos}

El sistema se organiza en \textbf{seis módulos}, presentados en la Tabla~\ref{tab:modulos}.

{\estilotabla
\begin{longtable}{L{0.28\textwidth} L{0.66\textwidth}}

\caption{Módulos funcionales del sistema}
\label{tab:modulos}
\\

\toprule
\textbf{Módulo} & \textbf{Descripción} \\
\midrule
\endfirsthead

\toprule
\textbf{Módulo} & \textbf{Descripción} \\
\midrule
\endhead

\textbf{M1. Administrar maestros} &
Configuración global del depósito: usuarios ilimitados con roles y permisos, clientes,
trabajadores, cuadrillas y tipos de bulto para parametrizar las operaciones. \\

\midrule

\textbf{M2. Ingesta de información de carga} &
Alimentación del sistema con el detalle de lo que viene dentro del contenedor, por tres vías
alternativas (ver Sección~\ref{sec:modalidades}). \\

\midrule

\textbf{M3. Programar faenas} &
Planificación de la agenda de desconsolidados: asignación de contenedor, fecha, franja horaria,
cuadrilla y jefe tarjador responsable. \\

\midrule

\textbf{M4. Ejecutar operaciones} &
Registro paso a paso de la faena en terreno desde la aplicación móvil, con captura obligatoria de
evidencia fotográfica en cada etapa crítica. \\

\midrule

\textbf{M5. Gestión de reportes} &
Compilación automática de datos y fotografías en un reporte PDF, y flujo de revisión y aprobación
en escritorio antes de su publicación. \\

\midrule

\textbf{M6. Publicar y enviar} &
Envío del reporte aprobado al cliente por correo electrónico y consulta histórica de faenas por
evento o por cliente. \\

\bottomrule

\end{longtable}
}

\section{Modalidades de ingreso de la información de carga}
\label{sec:modalidades}

Antes de abrir un contenedor, el sistema necesita saber qué hay dentro. BeeTracer ofrece
\textbf{tres vías}, en orden de preferencia:

\begin{enumerate}
    \item \textbf{Vía manifiesto electrónico (XML)} ---\este{modalidad ideal}---. El sistema se
          conecta al \este{web service} del Servicio Nacional de Aduanas, descarga el XML del
          manifiesto que envió la nave y precarga automáticamente contenedores, líneas de carga,
          consignatarios y cantidades declaradas. Permite conocer el contenido \textbf{antes de que
          el barco recale en el puerto}.
    \item \textbf{Integración directa (API)}: el sistema del propio cliente se conecta vía API con
          BeeTracer y le transfiere la información de la carga.
    \item \textbf{Tarja ciega} ---\este{modalidad de contingencia}---. Cuando no se recibió
          información previa, la aplicación móvil arranca «en blanco» y el operario ingresa
          manualmente la descripción de cada lote, leyendo las viñetas y facturas adheridas a las
          propias cajas a medida que las descarga.
\end{enumerate}

\section{El flujo operativo (\este{core} del sistema)}
\label{sec:flujo}

\textbf{Fase A --- Ingreso de información (Web).} Se carga el detalle de la carga por alguna de las
tres vías anteriores.

\noindent \textbf{Fase B --- Planificación (Web).} El Jefe de faena revisa los contenedores por arribar y
programa los desconsolidados, asignando franja horaria (p. ej. de 10:00 a 12:00) y una cuadrilla de
cuatro a cinco personas. A uno de sus integrantes se le asigna la responsabilidad de operar la
aplicación móvil: el \textbf{jefe tarjador}.

\noindent \textbf{Fase C --- Ejecución en terreno (App móvil).} El jefe tarjador sigue un flujo estricto que
la aplicación impone, detallado en la Tabla~\ref{tab:etapasterreno}.

{\estilotabla
\begin{longtable}{L{0.28\textwidth} L{0.66\textwidth}}

\caption{Etapas de la ejecución en terreno y evidencia exigida}
\label{tab:etapasterreno}
\\

\toprule
\textbf{Etapa} & \textbf{Registro exigido por el sistema} \\
\midrule
\endfirsthead

\toprule
\textbf{Etapa} & \textbf{Registro exigido por el sistema} \\
\midrule
\endhead

\textbf{1. Antes de abrir} &
Fotografía de las puertas, del candado y del sello de seguridad, más la \textbf{digitación
obligatoria del número de sello}. \\

\midrule

\textbf{2. Apertura} &
Fotografías panorámicas de la mercancía con el contenedor recién abierto, \textbf{antes de ser
manipulada}. \\

\midrule

\textbf{3. Separación} &
A medida que la cuadrilla descarga, el operario agrupa la mercancía por cliente y fotografía
\textbf{cada lote individualmente} (p. ej., las 200 cajas de zapatos). \\

\midrule

\textbf{4. Control de incidencias} &
Ante una diferencia de cantidad (faltante o sobrante) o una caja dañada, se separa la unidad
afectada y se le toman \textbf{fotografías exclusivas}. \\

\midrule

\textbf{5. Cierre} &
Fotografía del contenedor \textbf{completamente vacío}, para demostrar que no quedó carga
dentro. \\

\bottomrule

\end{longtable}
}

\noindent \textbf{Fase D --- Control y envío del reporte (Web).} Al cerrar la faena, el sistema compila
automáticamente todos los datos y fotografías en un \textbf{reporte PDF} con encabezado (fecha,
hora, cuadrilla a cargo) y desglose de fotografías y detalle por cliente, \textbf{destacando en
rojo} la mercancía reportada como dañada. El Supervisor revisa este borrador en la plataforma web;
si detecta una fotografía defectuosa ---borrosa, apuntando al cielo, o que muestra un detalle menor
que el terminal prefiere resolver internamente--- puede solicitar su eliminación y recaptura. Una
vez aprobado, el PDF completo se envía al cliente con un solo botón.

\begin{nota}
\textbf{Decisión de diseño relevante.} Originalmente el cliente accedía a un enlace y veía las
fotografías \textbf{en vivo}, a medida que se subían. Esta funcionalidad se sustituyó por el flujo
de revisión en escritorio, precisamente por los problemas de calidad y de exposición comercial
descritos en la Sección~\ref{sec:restricciones}.
\end{nota}

\section{Soluciones técnicas al entorno adverso}
\label{sec:soluciones}

{\estilotabla
\begin{longtable}{L{0.30\textwidth} L{0.64\textwidth}}

\caption{Soluciones técnicas adoptadas frente a las restricciones del entorno}
\label{tab:solucionestecnicas}
\\

\toprule
\textbf{Desafío} & \textbf{Solución implementada} \\
\midrule
\endfirsthead

\toprule
\textbf{Desafío} & \textbf{Solución implementada} \\
\midrule
\endhead

Wi-Fi bloqueado por los contenedores metálicos &
Envío \textbf{incremental}: la aplicación sube cada fotografía \textbf{una por una, inmediatamente
después de ser tomada}, transfiriendo pocos \este{bytes} a la vez, en lugar de enviar el lote
completo al final (lo que provocaba caídas). \\

\midrule

Peso de las imágenes &
\textbf{Compresión estratégica} en el dispositivo, buscando el punto de equilibrio entre bajo peso
y legibilidad suficiente para servir como evidencia. \\

\midrule

Ausencia total de señal &
\textbf{Modo offline 100\,\%}: la aplicación mantiene una base de datos local en el dispositivo y
sincroniza al recuperar conectividad o al llegar a una oficina con internet. \\

\bottomrule

\end{longtable}
}

\begin{nota}
En la práctica, la solución definitiva en la mayoría de los casos fue que los propios terminales
mejoraran la iluminación y la infraestructura Wi-Fi de sus patios.
\end{nota}

\section{Valor para el negocio}
\label{sec:valor}

Más allá de la eficiencia operativa, la función crítica del sistema es de \textbf{protección
legal}. Cuando el operario abre el contenedor y fotografía de inmediato el estado de la mercancía,
el depósito puede demostrar si la carga \textbf{ya venía dañada desde el barco} o si el daño
ocurrió accidentalmente durante sus propias maniobras de desconsolidado. Esto permite
\textbf{determinar responsabilidades ante los seguros de forma rápida y objetiva}, que es
precisamente lo que el proceso manual no permitía.
